% JAC1964-20261121. /root/programme_audit; GPT-6 Astra Pro.
% Additional established input: Kimura finite-dimensionality/conservativity.

\subsection*{JAC-P13. The canonical invariant complement}

\begin{proposition}[JAC-P13]\label{jac:P13}
Retain the Fermat notation of AR-8 and JAC-P11. Put
$P=(h^3(J),\pi_P)$, and let $e'$ be AR-8's self-transpose
Picard-summand projector with $e'p_G=p_Ge'=0$.
Then, as actual rational Chow projectors on $h^3(J)$,
\[
 \pi_P-e'=p_G\pi_P=p_G.                               \tag{JAC.P13.1}
\]
In particular the complementary motive is exactly
$N=h^3(J)^G$, independent of the auxiliary inclusion used in AR-2,
and
\[
 \operatorname{CH}^2(N)
       =\operatorname{CH}^2_{(1)}(J)_{\mathbf Q}^{G}.    \tag{JAC.P13.2}
\]
The operation $T$ of JAC-P11 consequently preserves all of
$\operatorname{CH}^2(N)$ and makes it a $\mathbf Q(i)$-module.
Its height is the real part of the complex-valued Hermitian form
\[
 \mathscr H(z,w)=h_K(z,w)+i\,h_K(z,Tw),                 \tag{JAC.P13.3}
\]
where $i$ acts through $T$ and the form is linear in its first
variable. Positivity of $\mathscr H$ on the whole module remains
the full complementary height problem.
\end{proposition}

\begin{proof}
We explicitly add the established Kimura input: a direct summand
of an abelian Chow motive is finite-dimensional; a finite-dimensional
Chow motive with zero Weil realization is zero.
These are Kimura, \emph{Chow groups are finite dimensional, in some
sense}, Math.\ Ann.\ 331 (2005), Example~9.1 and Corollary~7.3.
The exact application is also recalled and used in the directly
read proof of Theorem~A of Laga--Shnidman,
\emph{Vanishing criteria for Ceresa cycles}, author PDF p.~8.
The original Kimura article was not directly read in this run;
the cited standard consequence is carried through that primary
research application.

Set $q=\pi_P-e'-p_G\pi_P$. All three relevant projectors commute:
$e'$ is supported on $P$, $p_G$ commutes with $\pi_P$, and
$e'p_G=p_Ge'=0$. Therefore $q^2=q$.
AR-8's explicit character calculation shows that $p_GH^3(J)$ is
exactly the two-dimensional sum of the extreme Hodge types.
AR-2 identifies $H(e')$ with the remaining maximal level-one
part of $H(P)$. Hence $H(q)=0$.
The motive defined by $q$ is a direct summand of $h(J)$.
Kimura's consequence forces it, and thus the Chow projector $q$,
to be zero. This proves the first equality in (JAC.P13.1).
Also $p_G(1-\pi_P)$ is an idempotent on $h^3(J)$ with zero
realization, because the two extreme types are primitive.
The same argument gives $p_G=p_G\pi_P$.

These arguments can be applied after passing to an algebraic
closure. The resulting rational-equivalence identities descend:
a rational equivalence over the algebraic closure is defined over
a finite extension, and pushforward multiplies its descended class
by the extension degree. Thus no descent ambiguity is hidden here.
The Chow--Kunneth grading gives (JAC.P13.2).
Centrality makes $T$ commute with $p_G$, and JAC-P11 gives
$T^2=-1$ and height isometry. Therefore
$h_K(Tz,w)=-h_K(z,Tw)$.
This verifies conjugate symmetry and
$\mathscr H(Tz,w)=i\mathscr H(z,w)$,
$\mathscr H(z,Tw)=-i\mathscr H(z,w)$.
Finally $\mathscr H(z,z)=h_K(z,z)$, so their positivity questions
are identical. Nothing here bounds the dimension of this
$\mathbf Q(i)$-module.
\end{proof}


\subsection*{JAC-P14. The entire abelian-endomorphism orbit}

\begin{proposition}[JAC-P14]\label{jac:P14}
Choose an isogeny $J\sim E^3$ over $K$, with
$\operatorname{End}^0(E)=F=\mathbf Q(i)$, identifying
$\operatorname{End}^0(J)$ with $M_3(F)$. Choose the sign of $j$
in JAC-P11 so that it corresponds to the scalar matrix $iI_3$.
For $b\in\operatorname{End}^0(J)$ choose $d>0$ with $db$ an
actual endomorphism and define its normalized middle pullback by
\[
 \rho_3(b)=d^{-3}(db)^*\bigm|_{h^3(J)}.
\]
Then $\rho_3(b)$ preserves the exact Chow summand $N$ of JAC-P13.
If $\det_F(b)=u+vi$, with $u,v\in\mathbf Q$, its action is
\[
 \rho_3(b)\bigm|_N=u\,1_N-v\,T.                       \tag{JAC.P14.1}
\]
In particular
\[
 \operatorname{span}_{\mathbf Q}
   \{\rho_3(b)\eta:b\in\operatorname{End}^0(J)\}
        =\mathbf Q\eta\oplus\mathbf Q T\eta.           \tag{JAC.P14.2}
\]
The equality is at the level of rational Chow groups, not merely
after realization. For any $b,c\in\operatorname{End}^0(J)$,
\[
 h_K(\rho_3(b)\eta,\rho_3(c)\eta)
   =\operatorname{Re}
       \bigl(\det_F(b)\overline{\det_F(c)}\bigr)
      h_K(\eta,\eta).                                \tag{JAC.P14.3}
\]
Thus the full matrix endomorphism algebra does not enlarge the
positive CM plane.
\end{proposition}

\begin{proof}
We use the functorial Chow-motive power formula
$h^r(A)=\operatorname{Sym}^r h^1(A)$, with the usual commutativity
constraint on Chow motives. Its realization is the exterior power
of the underlying degree-one cohomology, because realization is
graded-commutative. The exact original reference is Kunnemann,
\emph{On the Chow motive of an abelian scheme}, in \emph{Motives},
Proc.\ Sympos.\ Pure Math.\ 55(1) (1994), 189--205,
Theorem~3.3.1. We directly checked the formulation in
Ancona--Enright-Ward--Huber, \emph{On the motive of a commutative
algebraic group}, Documenta Math.\ 20 (2015), 807--858,
Theorem~4.2.3(3), printed p.~825: the canonical isomorphism respects
the grading. Their Remark~4.2.1, printed p.~824, explains both
transposition to the contravariant convention used here and
Kunnemann's exterior-power notation for this symmetric-power functor.
The scalar normalization in $\rho_3$ is exactly the homogeneous
degree-three normalization in this functor; it is independent of $d$.

Extend the coefficient field of motives to $F$. The CM idempotents
split
\[
 h^1(E)_F=V_+\oplus V_-,
 \qquad j|_{V_+}=i,\quad j|_{V_-}=-i.
\]
Each $V_\pm$ has one-dimensional odd realization.
The motive $\operatorname{Sym}^2 V_\pm$ has zero realization and
is a summand of a tensor power of an abelian motive. By the
Kimura consequence explicitly invoked in JAC-P13 it is zero.
This consequence also holds after the coefficient extension to $F$:
an $F$-coefficient summand is a rational summand after restriction
of scalars, so its finite-dimensionality and the realization-zero
criterion follow from the rational statement.
In particular the interchange map on $V_\pm\otimes V_\pm$ is
exactly $-1$ in the category of Chow motives.

Now
\[
 h^1(J)_F\simeq V_+^{\oplus3}\oplus V_-^{\oplus3}.
\]
In its third symmetric power consider the two extreme summands
\[
 D_+=\operatorname{Sym}^3(V_+^{\oplus3}),\qquad
 D_-=\operatorname{Sym}^3(V_-^{\oplus3}).               \tag{JAC.P14.4}
\]
Since $\operatorname{Sym}^2V_\pm=0$, the direct-sum formula for
symmetric powers identifies $D_\pm$ with $V_\pm^{\otimes3}$:
the only surviving multiplicity pattern is $(1,1,1)$.
The interchange sign just proved shows that a matrix
$b\in M_3(F)$ acts on $D_+$ by $\det_F(b)$ and on $D_-$ by
$\overline{\det_F(b)}$. Indeed expansion of the action on its
three tensor factors is precisely the alternating permutation sum
defining the determinant. This calculation is in Chow motives.

It remains to identify these summands with the particular $N$,
rather than only with a motive having the same realization.
The sum of the projectors onto $D_+$ and $D_-$ is defined over
$\mathbf Q$. For an explicit description, apply $\rho_3(1+j)$.
Its eigenvalues on the two extreme summands are $-2\pm2i$,
while on the two mixed summands they are $2\pm2i$.
The polynomials
\[
 X^2+4X+8,\qquad X^2-4X+8
\]
are coprime over $\mathbf Q$, so their Chinese-remainder
idempotent defines an exact rational Chow projector $p_{\rm ext}$.
The preceding tensor decomposition verifies the polynomial
identity needed for this idempotent exactly, not only in cohomology.
Because $1+j$ is central, $p_{\rm ext}$ commutes with all
$\rho_3(b)$, and in particular with $p_G$.
The two commuting idempotents $p_{\rm ext}$ and $p_G$ have
identical realizations. Therefore
$p_{\rm ext}(1-p_G)$ and $p_G(1-p_{\rm ext})$ are
realization-zero idempotents. Kimura makes both zero.
JAC-P13 now gives $N=(h^3(J),p_G)=(h^3(J),p_{\rm ext})$.

On $D_+$ the operator $T=\rho_3(j)$ is multiplication by
$i^3=-i$, and on $D_-$ it is multiplication by $(-i)^3=i$.
Consequently $u-vT$ has eigenactions $u+vi$ and $u-vi$,
exactly the determinant actions computed above. This proves
(JAC.P14.1) after extension of coefficients to $F$ and hence over
$\mathbf Q$, since extension of a rational coefficient vector
space to $F$ is injective.
The identity endomorphism and $j$ yield $\eta$ and $T\eta$,
while (JAC.P14.1) puts every other image in their span.
This proves (JAC.P14.2). Finally JAC-P11 gives the diagonal
height matrix with equal entries and zero mixed entry, which
implies (JAC.P14.3).
\end{proof}

\paragraph{Interpretation of JAC-P13--P14.}
The remaining question is positivity on an actual
$\mathbf Q(i)$-module of primitive cycles in the exact motive
$h^3(J)^G$. The positive vector $\eta$ generates precisely the
CM plane under all abelian endomorphisms. Further progress on
the entire group therefore needs information on additional
$\mathbf Q(i)$-independent cycles, or a theorem proving there
are none; applying more abelian endomorphisms to $\eta$ cannot
supply that information.


